% !TeX root = ../../main.tex

Bei der Implementierung wurden verschiedene Entscheidungen getroffen, um den Versuchsaufbau möglichst einfach und
reproduzierbar zu halten.
Gleichzeitig sollte die Implementierung möglichst nah an den tatsächlich untersuchten Datenbank- und
Retry-Mechanismen bleiben.
Die wichtigsten Entscheidungen werden im Folgenden kurz begründet.

Für die Umsetzung wurde Python verwendet, da sich damit sowohl die Kommunikation mit PostgreSQL als auch die
parallele Ausführung der Transaktionen und die spätere Analyse der Versuchsdaten umsetzen lassen.
Für die Datenbankanbindung wird \texttt{psycopg2} eingesetzt, damit ein direkter Zugriff auf PostgreSQL
möglich ist.
Dadurch können SQL-Anweisungen und Transaktionen gezielt gesteuert und auftretende Datenbankfehler direkt
verarbeitet werden.
Zudem handelt es sich bei \texttt{psycopg2} um eine etablierte und weit verbreitete Bibliothek für die Anbindung
von PostgreSQL aus Python.
Als Datenbanksystem wird PostgreSQL verwendet, da die für die Experimente benötigten Transaktionsmechanismen und
Sperrverfahren unterstützt werden.
Insbesondere die Erkennung von Fehlern und die Rückgabe entsprechender SQLSTATE-Codes sind für die Untersuchung der
Fehlerklassen relevant.
Die Datenbank wird lokal betrieben, um zusätzliche Einflüsse durch Netzwerkverbindungen oder externe Systeme zu
vermeiden.
Dadurch kann die Testumgebung zudem einfacher konfiguriert und kontrolliert werden.

Für die Speicherung der Versuchsergebnisse wurde das JSONL-Format gewählt.
Jeder Transaktionsversuch kann dabei als eigener Datensatz gespeichert werden, wodurch auch bei einer großen Anzahl
an Messwerten eine einfache und strukturierte Speicherung möglich ist.
Die erzeugten Logdateien können anschließend direkt mit \texttt{pandas} eingelesen und weiterverarbeitet werden.
\texttt{pandas} wird für die Gruppierung der Versuchsdaten sowie die Berechnung der untersuchten Kenngrößen
eingesetzt.

Die Fehler werden gezielt erzeugt, um für jedes Experiment eine definierte Fehlerklasse untersuchen zu können.
Eine dynamische Klassifikation anhand der SQLSTATE-Codes ist daher für die Durchführung der Experimente nicht
notwendig.
Die Codes werden dennoch protokolliert, um die tatsächlich aufgetretenen Fehler überprüfen und die Ergebnisse
entsprechend nachvollziehen zu können.

Für die Untersuchung des Retry-Verhaltens wurden die Strategien \emph{Immediate Retry}, \emph{Retry after Delay} und
\emph{Retry with Backoff + Jitter} gewählt.
Diese decken unterschiedliche Ansätze der Wiederholungsversuche ab.
Während beim \emph{Immediate Retry} unmittelbar ein neuer Versuch gestartet wird, verwendet \emph{Retry after Delay}
eine konstante Wartezeit.
\emph{Retry with Backoff + Jitter} erweitert diesen Ansatz um einen exponentiell steigenden maximalen Delay und
eine zufällige Komponente.
Dadurch können die Auswirkungen verschiedener Verzögerungsstrategien auf die untersuchten Kenngrößen miteinander
verglichen werden.

Die parallele Ausführung der Clients erfolgt über den \texttt{ThreadPoolExecutor}.
Dadurch können mehrere Transaktionen gleichzeitig ausgeführt und die für die Experimente benötigte Concurrency
realisiert werden.
Der nächste Versuch wird erst gestartet, nachdem der vorherige fehlgeschlagen ist und gegebenenfalls der
entsprechende Delay abgewartet wurde.
